% ---------------------------------------------------------------
% Paper A -- Sections 3 and 4
% Requires: booktabs, amsmath, siunitx (optional), url
% ---------------------------------------------------------------

\section{Method}
\label{sec:method}

\subsection{Overview}

We evaluate carbon-aware placement of non-preemptible virtual machines across
geographically distributed datacentres. A workload of independent VM requests
arrives over a 720\,h horizon. Each request must be assigned to one region and
one start hour, subject to a per-region concurrency cap and a per-request
deadline. Once started, a VM runs to completion in the region where it was
placed; migration and preemption are outside the scope of this study, following
the assumptions of Zanotto et al.~\cite{zanotto2025usercentric}.

The evaluation is implemented on CloudSim Plus 8.5.7~\cite{silvafilho2017cloudsimplus}
under Java~17. Placement
decisions are taken by a planner operating on an occupancy ledger, and the
resulting schedule is then executed in the simulator so that host-level power
draw, and hence emissions, are obtained from simulated execution rather than
from the planner's own cost model. This separation matters: the planner
minimises a carbon-intensity objective, whereas the reported emissions depend
additionally on how many physical hosts the schedule activates. Sections
\ref{sec:accounting} and \ref{sec:validity} describe the accounting and the
checks that connect the two.

\subsection{System model}

Let $J$ denote the set of eligible regions and $T = \{0,\dots,H-1\}$ the hours
of the horizon, with $H = 720$. Each request $r$ is a tuple
$(p_r, m_r, d_r, a_r, \delta_r)$ giving its core count, memory, duration in
hours, arrival hour and deadline hour. A placement assigns $r$ to a region
$j_r \in J$ and a start hour $s_r \ge a_r$. The request meets its deadline when
$s_r + d_r \le \delta_r$.

Regional capacity is expressed as an occupancy matrix
$\mathrm{alloc}[j][t]$ counting VMs running in region $j$ at hour $t$, subject
to
%
\begin{equation}
\mathrm{alloc}[j][t] \;\le\; C_j
\qquad \forall j \in J,\; \forall t \in T ,
\label{eq:cap}
\end{equation}
%
where $C_j$ is the per-region concurrency cap. Constraint~\eqref{eq:cap}
represents the finite capacity a tenant can command in any single region and is
the parameter whose effect this paper isolates.

Each region hosts one datacentre containing an identical fleet of homogeneous
hosts. The host is a Dell PowerEdge XR8620T with an Intel Xeon Gold 6433N
(32 cores at 2.00\,GHz) and 256\,GB of memory. This is the platform modelled by
Zanotto et al.~\cite{zanotto2025usercentric}, who likewise draw its power profile from SPECpower and likewise
evaluate on CloudSim Plus against a carbon-agnostic round-robin baseline;
adopting the same host removes hardware as a source of difference between their
reported figures and ours. Each core is modelled as a processing element of
2000\,MIPS.
A VM requests $p_r$ processing elements and $m_r$ memory, and executes a single
cloudlet of length $2000 \cdot d_r \cdot 3600$ MI at full CPU utilisation, so
that its runtime equals $d_r$ hours.

\subsection{Host power model}
\label{sec:power}

Host power is taken from the published SPECpower\_ssj2008 result for the same
machine \cite{spec2024xr8620t} (Dell Inc., PowerEdge XR8620T with Xeon Gold 6433N,
tested 15 July 2024, 6\,564 overall ssj\_ops/watt). The sensitivity of simulated
energy figures to the power model adopted is well documented
\cite{ismail2020serverpower}. Because the simulated host matches the
benchmarked one in core count, clock and memory, the measured curve is used
without rescaling. Table~\ref{tab:power} reports it.

\begin{table}[t]
\centering
\caption{Measured host power at decile load points (SPECpower\_ssj2008).}
\label{tab:power}
\begin{tabular}{lrrrrrrrrrrr}
\toprule
Load (\%) & 0 & 10 & 20 & 30 & 40 & 50 & 60 & 70 & 80 & 90 & 100 \\
Power (W) & 230 & 244 & 260 & 273 & 301 & 340 & 372 & 385 & 400 & 416 & 420 \\
\bottomrule
\end{tabular}
\end{table}

Power at intermediate utilisations is obtained by linear interpolation between
adjacent measured points, since SPEC publishes no finer resolution. Two
properties of this curve affect the results. First, the machine draws 230\,W at
idle against a 420\,W peak, an idle fraction of 55\,\%; the range a scheduler
can influence is therefore 190\,W per host. Second, the curve is not linear: it
is nearly flat between 0\,\% and 30\,\% load, rises steeply between 30\,\% and
60\,\%, and flattens again above 70\,\%. A linear model fitted only to the idle
and peak points misprices both ends of the range traversed by the capacity
sweep. Section~\ref{sec:robustness-design} describes how sensitivity to this
choice is tested.

\subsection{Carbon accounting}
\label{sec:accounting}

Emissions are obtained by integrating host power against regional carbon
intensity. Sampling occurs on every simulation clock tick, with a datacentre
scheduling interval of 300\,s, and each interval is charged at the carbon
intensity of the hour in which it begins. For a host $h$ in region $j$ with
utilisation $u_h(t)$, the instantaneous rate is
$P(u_h(t)) \cdot \mathrm{CI}_j(t)$, where $P$ is the interpolated curve of
Table~\ref{tab:power} and $\mathrm{CI}_j(t)$ the regional carbon intensity in
gCO\textsubscript{2}eq/kWh.

We report emissions on three bases, because the choice determines what a
comparison can detect.

\begin{align}
E_{\mathrm{total}} &= \sum_{j \in J} \sum_{h \in \mathcal{H}_j} \int P(u_h(t))\,\mathrm{CI}_j(t)\,dt ,
\label{eq:total}\\[2pt]
E_{\mathrm{dyn}}   &= \sum_{j \in J} \sum_{h \in \mathcal{H}_j} \int \bigl[P(u_h(t)) - P(0)\bigr]\,\mathrm{CI}_j(t)\,dt ,
\label{eq:dyn}\\[2pt]
E_{\mathrm{attr}}  &= \sum_{j \in J} \sum_{h \in \mathcal{H}_j} \int \mathbb{1}\{u_h(t) > 0\}\, P(u_h(t))\,\mathrm{CI}_j(t)\,dt .
\label{eq:attr}
\end{align}

Equation~\eqref{eq:total} is the provider's meter: every provisioned host draws
power for the whole measurement window whether or not it carries load. Because
the same fleet exists in the same regions under every policy, the idle
component of $E_{\mathrm{total}}$ is invariant to placement. In our
configuration this component dominates, and two policies with entirely
different schedules report totals differing by less than 2\,\%. A study
reporting only $E_{\mathrm{total}}$ cannot distinguish the policies it compares.

Equation~\eqref{eq:dyn} counts only power above idle. It is sensitive to
placement and is unaffected by how many hosts a schedule activates, which makes
it the appropriate basis when the quantity of interest is the carbon cost of
the work itself.

Equation~\eqref{eq:attr} charges the full power of hosts that carry at least one
VM and nothing for hosts that carry none. This is the consumer-side basis: a
tenant is accountable for the machines its workload occupies, not for the
provider's idle capacity. It is sensitive both to regional placement and to
consolidation, and it is the basis we report as primary, with
$E_{\mathrm{dyn}}$ given alongside so that the consolidation component of any
reduction remains visible.

All three are normalised by admitted VM-hours to permit comparison across
configurations.

\subsection{Scheduling policies}

Four policies are compared. All see the same carbon trace and the same
occupancy ledger, so differences arise from the decision rule alone.

\emph{Round robin} places each request at its arrival hour in the next region
with available capacity, ignoring carbon intensity. It is the carbon-agnostic
baseline and is included so that the share of any reported reduction
attributable to baseline weakness remains visible.

\emph{Space} places each request at its arrival hour in the eligible region
with the lowest mean carbon intensity over the request's execution window. It
shifts across regions but never defers.

\emph{WaitAwhile}, after Wiesner et al.~\cite{wiesner2021letswait}, holds each request in its origin
region and defers it to the cleanest feasible start hour within its deadline.
The origin region is a parameter of the policy, denoted
$\mathrm{WaitAwhile}(j_0)$; Section~\ref{sec:design-factors} explains why it is
swept rather than fixed.

\emph{Space+time} searches jointly over regions and start hours within the
deadline. Zanotto et al.~\cite{zanotto2025usercentric} solve the corresponding problem as a
binary integer program applied to one request at a time against a greedily updated occupancy
matrix. For a single request the feasible set has $|J| \times |T|$ members, so
exhaustive enumeration returns the same optimum without a solver dependency.

Every policy shares a common fallback: when no feasible slot exists within the
deadline, the request is placed at the earliest feasible slot anywhere,
accepting a deadline violation. The fallback searches all regions, including for
WaitAwhile, since a policy that declines to relocate once its deadline is
already unreachable would be constrained beyond what its formulation states.
Requests that cannot be placed even beyond their deadline cause the run to
fail rather than to be silently dropped.

\subsection{Placement validity}
\label{sec:validity}

Two properties of CloudSim Plus can silently break the correspondence between
the planned schedule and the executed one. A datacentre that cannot
accommodate a VM causes the broker to reroute it to another datacentre, so that
the executed placement differs from the planned region; and the broker's
datacentre mapper is applied per batch of waiting VMs rather than per VM, so
that a single broker cannot express per-VM regional assignment at all. We
therefore instantiate one broker per region, each pinned to its own datacentre,
which makes regional misplacement impossible by construction.

After each run we verify, for every VM, that it was created in the datacentre
its planner selected and that its cloudlet completed within the measurement
window. Placement is recovered from each broker's list of created VMs rather
than from the VM's own creation flag, which reads false once a VM has been
destroyed at the end of its execution. Cells failing this check are excluded
from analysis rather than reported, since their emissions do not describe the
schedule under test.

\section{Experimental design}
\label{sec:design}

\subsection{Carbon intensity traces}

Carbon intensity is taken from CarbonCast~\cite{maji2022carboncast}, which publishes
hourly lifecycle emission factors for thirteen grid regions across North America, Europe and
Australia for 2020--2021. We use it in preference to commercial sources because
the data may be redistributed, which allows the experiment to be reproduced
from a public archive, unlike the metered access of commercial providers
\cite{electricitymaps}, and because it also publishes paired actual and
forecast series, which the companion study requires.

All thirteen regional files were verified before use: each contains 17\,544
hourly observations with no missing hours, no duplicate timestamps and no
non-positive values. Timestamps are not written in a uniform format across
regions, and files are aligned on a common UTC window rather than on each
file's own first row, since indexing each region independently would offset
the regions against one another.

Zanotto et al.~\cite{zanotto2025usercentric} site their clusters at Azure cloud regions. We instead use grid
regions for which independently published hourly emission factors exist, since
the object of study is the effect of region-set composition and that requires
regional carbon intensity to be a measured input rather than a modelling
choice. Two region sets are used. The primary set comprises seven United States
independent system operators; the second is a five-region European set retained
as a control. Table~\ref{tab:regions} reports their characteristics over the
720\,h window beginning 1 January 2020. The column headed \emph{diurnal} gives
the spread of hour-of-day means as a percentage of the regional mean, which
isolates the repeatable daily cycle available to a deferral policy from
one-off excursions.

\begin{table}[t]
\centering
\caption{Regional characteristics over the 720\,h window from 1 January 2020.
Mean carbon intensity in gCO\textsubscript{2}eq/kWh; diurnal amplitude as a
percentage of the regional mean.}
\label{tab:regions}
\begin{tabular}{llrr}
\toprule
Set & Region & Mean CI & Diurnal (\%) \\
\midrule
\multirow{7}{*}{US}
 & BPAT  &  63 & 14.3 \\
 & NYISO & 209 & 25.6 \\
 & ISNE  & 297 & 14.5 \\
 & CISO  & 299 & 52.7 \\
 & PJM   & 341 &  7.8 \\
 & ERCO  & 354 & 24.0 \\
 & FPL   & 359 & 21.1 \\
\midrule
\multirow{5}{*}{EU}
 & SE & 39 &  1.7 \\
 & ES & 183 & 13.7 \\
 & DE & 323 & 15.3 \\
 & NL & 483 &  8.4 \\
 & PL & 663 &  6.1 \\
\bottomrule
\end{tabular}
\end{table}

The two sets differ in a way that is central to
Section~\ref{sec:design-factors}. In the US set the cleanest region is 3.3
times cleaner than the next cleanest, and the full spread is 5.7-fold. In the
EU set the cleanest region is 4.7 times cleaner than the next, and the spread is
17.1-fold. The EU set therefore admits a dominant optimum, and is retained
precisely so that the consequences of that dominance can be measured.

\subsection{Workload}

VM requests are sampled from the Azure Resource Central
trace~\cite{cortez2017resourcecentral,azurepublicdataset}, which records 2\,695\,549 virtual machines with creation and
deletion timestamps, core count, memory and a workload category assigned by the
provider. Timestamps run to 2\,591\,400\,s, exactly the 720\,h horizon
simulated here, so no rescaling of arrival times is required.

Three filters are applied. Virtual machines still running at the end of the
trace (224\,249 records) are right-censored and are excluded, because treating
the truncation point as a completion would systematically shorten the
long-running VMs that deferral is intended to target. Virtual machines shorter
than two hours are excluded, since a VM shorter than the scheduling granularity
cannot meaningfully be deferred. Sampling is by reservoir rather than by taking
a prefix of the file, since the trace is not shuffled.

A sample of 1\,000 requests drawn under these filters yields approximately
33\,000 VM-hours over the horizon, a mean concurrency of about 46 VMs and a
mean VM lifetime of 34\,h. The provider's categories are unevenly distributed:
2\,446\,819 records are labelled \emph{Unknown}, 159\,497
\emph{Delay-insensitive} and 78\,309 \emph{Interactive}. A separate sample
restricted to the delay-insensitive class is used as a robustness condition.
That class has a mean lifetime of 340.7\,h, an order of magnitude longer than
the unrestricted mixture, so the sample size is reduced to 100 requests to hold
total demand, and hence capacity utilisation, constant across the two
conditions.

\subsection{Factors and levels}
\label{sec:design-factors}

Four factors are varied.

\emph{Capacity.} Total concurrent slots take four levels, distributed evenly
across the regions of the set. For the US set these are 56, 84, 140 and 350
slots, corresponding to capacity utilisations of 82\,\%, 55\,\%, 33\,\% and
13\,\% respectively; the EU levels are chosen to match those utilisations. The
sweep is defined on total slots rather than on a per-region cap because a
per-region cap makes total capacity depend on the size of the region set, so
that a policy confined to one region would compete against policies commanding
$|J|$ times more hardware.

\emph{Deadline margin.} The slack allowed beyond a request's own duration takes
values of 6, 12, 24 and 48\,h.

\emph{Home region.} WaitAwhile is instantiated once per region of the set
rather than once per configuration. A single instantiation reports the
performance of one arbitrary origin, and in both region sets the conventional
choice of the first listed region selects the cleanest grid available, which
assigns the policy the optimum without stating that it has done so.

\emph{Region set.} The US and EU sets described above.

Each configuration is repeated across five independent workload samples for the
US set and three for the EU set, obtained by varying the reservoir sampling
seed with all other inputs held fixed. Reported spreads are standard deviations
across these repetitions and represent workload sampling variability.

\subsection{Controls}
\label{sec:robustness-design}

Several quantities are deliberately held constant so that the factors above are
not confounded.

The physical fleet is sized once from the largest capacity level in the sweep
and held fixed across all cells. Sizing each cell to its own cap would make
total emissions a function of the capacity axis through the quantity of
hardware provisioned rather than through scheduling. Fleet size is derived from
the largest VM in the workload rather than the mean, because the cap bounds the
number of concurrent VMs and not their size, and CloudSim Plus allocates from
individual hosts rather than from a pooled resource.

All policies are measured over an identical window of $H + 48$ hours. Without a
fixed window a policy that defers work would run for longer and accrue
additional idle energy, which would be read as higher emissions irrespective of
its placement decisions. The ledger extends a further 336\,h beyond the horizon
so that deadline-violating fallbacks always have a feasible slot.

Planning is performed against the ground-truth carbon trace throughout, so the
policies operate under perfect foresight. This is intentional. Forecast error
would interact with both capacity and deadline margin, and the object of this
study is the effect of those two parameters in isolation; the effect of
prediction quality is addressed separately. Carbon-intensity forecasts of the
kind required for that analysis are themselves an active subject
\cite{maji2022carboncast,maji2022dacf}.

Finally, the entire experiment is repeated under a linear power model with a
120\,W idle and 480\,W peak, an idle fraction of 25\,\% against the 55\,\% of
the measured curve, in order to establish whether the reported effects depend
on the power model adopted.
